% =============================================================================
% Manual P03 – Anemometría de Hilo Caliente (CTA/HWA)
% Técnicas de Medida – Ingeniería Aeronáutica
% =============================================================================
\documentclass[12pt,a4paper]{article}

\usepackage[utf8]{inputenc}
\usepackage[T1]{fontenc}
\usepackage[spanish]{babel}
\usepackage[top=2.5cm,bottom=2.5cm,left=3cm,right=2.5cm]{geometry}
\usepackage{amsmath,amssymb,amsfonts}
\usepackage{siunitx}
\sisetup{output-decimal-marker={,}, group-separator={.}}
\usepackage{graphicx,xcolor,float,caption,subcaption}
\usepackage{tikz}
\usetikzlibrary{arrows.meta,positioning,calc,shapes.geometric,
                decorations.markings,patterns,fit,backgrounds}
\usepackage{circuitikz}
\usepackage{booktabs,multirow,array}
\usepackage{listings}
\lstset{
  basicstyle=\ttfamily\small,
  keywordstyle=\color{blue!70!black}\bfseries,
  commentstyle=\color{green!50!black},
  stringstyle=\color{red!60!black},
  backgroundcolor=\color{gray!10},
  frame=single, breaklines=true,
  numbers=left, numberstyle=\tiny\color{gray}
}
\usepackage[colorlinks=true,linkcolor=blue!70!black,citecolor=green!50!black,
            urlcolor=blue!50!black]{hyperref}
\usepackage[backend=biber,style=ieee,sorting=none]{biblatex}
\addbibresource{references_P03.bib}
\usepackage{fancyhdr}
\pagestyle{fancy}
\fancyhf{}
\fancyhead[L]{\small Técnicas de Medida -- Ingeniería Aeronáutica}
\fancyhead[R]{\small P03: Anemometría de Hilo Caliente}
\fancyfoot[C]{\thepage}

\definecolor{aeroblue}{RGB}{0,51,102}
\definecolor{aerored}{RGB}{153,0,0}
\definecolor{aerogreen}{RGB}{0,102,51}
\definecolor{aeroorange}{RGB}{200,100,0}

% =============================================================================
\begin{document}

\begin{titlepage}
  \centering
  \vspace*{1cm}
  {\color{aeroblue}\rule{\linewidth}{2pt}}\\[0.4cm]
  {\Large\bfseries\color{aeroblue} TÉCNICAS DE MEDIDA EN INGENIERÍA AERONÁUTICA}\\[0.3cm]
  {\color{aeroblue}\rule{\linewidth}{1pt}}\\[1.5cm]
  {\huge\bfseries P03 -- Anemometría de\\[0.3cm]Hilo Caliente (CTA/HWA)}\\[1cm]
  {\large Manual de Laboratorio}\\[0.5cm]
  {\large\itshape Ingeniería Aeronáutica}\\[2cm]
  \begin{tikzpicture}[scale=0.9,>=Stealth]
    % Wire probe
    \draw[fill=gray!30,draw=black,thick] (0,-0.5) rectangle (0.4,2.5);
    \draw[fill=gray!50,draw=black,thick] (0,2.5) -- (0.2,2.7) -- (0.4,2.5);
    % Prongs
    \draw[black,very thick] (0.2,2.7) -- (-0.3,3.3);
    \draw[black,very thick] (0.2,2.7) -- (0.7,3.3);
    % Wire (the sensing element)
    \draw[aerored,ultra thick,line width=2pt] (-0.3,3.3) -- (0.7,3.3);
    \node[aerored,below,font=\small] at (0.2,3.15) {Hilo ($d\approx\SI{5}{\micro\m}$)};
    % Heat transfer arrows
    \foreach \x in {-0.5,-0.15,0.2,0.55,0.9} {
      \draw[->,blue!60,thick] (\x,3.8) -- (\x,3.35);
    }
    \node[blue!60,above,font=\small] at (0.2,3.85) {Flujo $U$};
    % CTA box
    \draw[fill=aeroblue!20,draw=aeroblue,very thick] (2.0,2.8) rectangle (4.5,3.8);
    \node[aeroblue,font=\small\bfseries] at (3.25,3.45) {CTA Bridge};
    \node[font=\tiny] at (3.25,3.1) {Dantec StreamLine};
    % Connection
    \draw[aeroblue,thick,->] (0.4,3.3) -- (2.0,3.3);
    \node[above,font=\tiny] at (1.2,3.35) {$R_w(T)$};
    % Output to DAQ
    \draw[aeroblue,thick,->] (4.5,3.3) -- (5.5,3.3)
          node[right,font=\small] {$E(t)$};
  \end{tikzpicture}\\[1.5cm]
  {\large\textbf{Departamento de Ingeniería Aeroespacial}}\\[0.2cm]
  {\normalsize Laboratorio de Técnicas de Medida}\\[0.5cm]
  {\normalsize \today}
  \vfill
  {\color{aeroblue}\rule{\linewidth}{2pt}}
\end{titlepage}

\tableofcontents
\newpage

% =============================================================================
\section{Introducción}
\label{sec:intro}

La \textbf{anemometría de hilo caliente} (Hot-Wire Anemometry, HWA) es una de
las técnicas de medida de velocidad más versátiles en flujos aerodinámicos.
Se basa en la transferencia de calor por convección entre un hilo metálico
(sensor) calentado eléctricamente y el flujo de gas que lo rodea.  Su tiempo
de respuesta extremadamente corto (hasta \SI{400}{\kHz}) la convierte en la
herramienta de referencia para la caracterización de flujos turbulentos,
espectros de energía, e intermitencia de transición laminar-turbulenta en
aeronáutica \cite{bruun1995}.

\subsection{Aplicaciones en ingeniería aeronáutica}

\begin{itemize}
  \item \textbf{Capa límite sobre alas}: perfil de velocidades $U(y)$,
        posición de transición laminar-turbulenta, grosor de desplazamiento
        $\delta^*$.
  \item \textbf{Estela de perfiles}: vorticidad, déficit de arrastre, espectros.
  \item \textbf{Entrada de motores turbofán}: distorsión de flujo, total
        pressure recovery.
  \item \textbf{Flujo en toberas y difusores}: perfiles turbulentos, intensidad
        de turbulencia $I_u$.
  \item \textbf{Caracterización de turbulencia en secciones de prueba de túneles}.
\end{itemize}

% =============================================================================
\section{Objetivos}
\label{sec:objetivos}

\begin{enumerate}
  \item Comprender el principio de transferencia de calor del hilo caliente y
        la ley de King.
  \item Calibrar el sensor HWA obteniendo la curva $E^2$ vs.\  $U^{1/n}$
        (ley de King linealizada).
  \item Determinar las constantes $A$, $B$ y $n$ de King para el sensor
        del laboratorio.
  \item Calcular estadísticos de turbulencia: $U_\mathrm{mean}$, $u'_\mathrm{rms}$,
        intensidad de turbulencia $I_u$, skewness y kurtosis.
  \item Obtener el espectro de densidad de potencia (PSD) de la señal de
        velocidad y comparar con la ley de Kolmogorov ($k^{-5/3}$).
\end{enumerate}

% =============================================================================
\section{Fundamento Teórico}
\label{sec:teoria}

\subsection{Transferencia de calor del hilo}

Un hilo cilíndrico de diámetro $d$ a temperatura $T_w$ en un flujo de
temperatura $T_a$ pierde calor a razón:
\begin{equation}
  Q_\mathrm{conv} = h\,A_s\,(T_w - T_a)
  \label{eq:calor_conv}
\end{equation}

con $A_s = \pi d L$ la superficie del hilo de longitud $L$ y $h$ el coeficiente
de convección.  En régimen estacionario, la potencia eléctrica disipada en el
hilo iguala la potencia convectiva:
\begin{equation}
  P_e = I_w^2\,R_w = Q_\mathrm{conv}
  \label{eq:balance_energia}
\end{equation}

El número de Nusselt para un cilindro en flujo cruzado se describe bien
mediante:
\begin{equation}
  Nu = \frac{hd}{k_f} = A_0 + B_0\,Re^n
  \label{eq:Nu_Re}
\end{equation}

donde $k_f$ es la conductividad térmica del fluido y $Re = Ud/\nu$.

\subsection{Ley de King}

Combinando las expresiones anteriores, la tensión $E$ a los bornes del puente
CTA satisface la \textbf{ley de King}:
\begin{equation}
  \boxed{E^2 = A + B\,U^n}
  \label{eq:king}
\end{equation}

con $A$, $B$ constantes del sensor y $n \approx 0.45$ (King original) o
$n = 0.5$ (ley de la raíz cuadrada).  En la práctica se ajusta $n$ como
parámetro libre, obteniéndose habitualmente $n \in [0.4, 0.5]$.

\paragraph{Linealización.}  Elevando ec.~\eqref{eq:king} a la potencia
$1/(2n)$ no es conveniente.  Se usa la forma linealizada en $E^2$ y
$U^{1/n}$:
\begin{equation}
  E^2 = A + B\,U^n
  \;\Longleftrightarrow\;
  E^2 = B\, (U^{1/n})^n + A
  \label{eq:king_lin}
\end{equation}

Definiendo $x = U^{1/n}$:  $E^2 = B\,x^n + A$, que es lineal en $x$ para
$n=1$ (no usual) o requiere ajuste no lineal en la forma general.

\subsection{Anemómetro a temperatura constante (CTA)}

El puente Wheatstone mantiene $R_w = \mathrm{cte}$ (y por tanto $T_w =
\mathrm{cte}$) ajustando la corriente $I_w$ mediante un amplificador de alta
ganancia.

\begin{figure}[H]
  \centering
  \begin{circuitikz}[scale=0.95, transform shape]
    % Wheatstone bridge
    % Top node
    \node (A) at (0,2) {};
    % Bottom node
    \node (B) at (0,-2) {};
    % Left node
    \node (L) at (-2,0) {};
    % Right node
    \node (R) at (2,0) {};

    % Bridge arms
    \draw (A) to[R,l=$R_1$] (L);
    \draw (A) to[R,l=$R_2$] (R);
    \draw (L) to[R,l=$R_3$,color=aerogreen] (B);
    % Hot wire arm (Rw)
    \draw (R) to[R,l_=$R_w(T_w)$,color=aerored] (B)
          node[midway,right,font=\tiny,aerored] {Sensor};
    % Power supply
    \draw (A) to[battery1,l=$V_s$] (0,3.5);
    \draw (0,-3.5) to[short] (B);
    \draw (0,3.5) to[short] (0,3.5);
    \draw (0,-2) to[short] (0,-3.5);

    % Differential amplifier between L and R
    \draw (L) -- (-2.5,0) -- (-2.5,0.5);
    \draw (R) -- (2.5,0) -- (2.5,0.5);
    % Op-amp
    \draw[fill=aeroblue!10] (-1.5,0.5) -- (1.5,0.5) -- (0,1.8) -- cycle;
    \node[font=\tiny,aeroblue] at (0,1.0) {Amp.};
    % Output
    \draw[aeroblue,thick,->] (0,1.8) -- (0,2.8)
          node[above,font=\small] {$E_\mathrm{out}(U)$};
    % Feedback to bridge
    \draw[aeroorange,dashed,->] (0,2.5) -- (3.5,2.5) -- (3.5,-1) -- (2.2,-0.2);
    \node[aeroorange,font=\tiny] at (3.8,1) {Feedback};
    % Labels
    \node[font=\tiny,left] at (-2.0,0) {$V_L$};
    \node[font=\tiny,right] at (2.0,0) {$V_R$};
    \node[font=\tiny] at (0,-0.3) {$\Delta V = V_R - V_L \to 0$};
    \draw[<->,gray,thin] (-3.5,-2) -- (-3.5,2)
          node[midway,left,font=\tiny] {Puente\\Wheatstone};
  \end{circuitikz}
  \caption{Puente Wheatstone del anemómetro CTA.  El amplificador de alta
           ganancia ajusta la corriente $I_w$ para mantener $R_w$ (y $T_w$)
           constante.  La tensión de salida $E_\mathrm{out}$ es proporcional
           a la corriente de calentamiento y codifica la velocidad del flujo.}
  \label{fig:CTA_circuit}
\end{figure}

\subsection{Resistencia del hilo en función de la temperatura}

La resistencia del hilo (típicamente wolframio) varía linealmente con $T$:
\begin{equation}
  R_w = R_0\,[1 + \alpha_0\,(T_w - T_0)]
  \label{eq:Rw_T}
\end{equation}

con $\alpha_0 \approx \SI{4.5e-3}{\per\kelvin}$ (W, a $T_0 = \SI{20}{\degreeCelsius}$).

La relación de sobrecalentamiento (overheat ratio):
\begin{equation}
  a = \frac{R_w - R_a}{R_a} = \alpha_0\,(T_w - T_a)
  \label{eq:overheat}
\end{equation}

Un valor típico es $a = 0.8$ ($T_w \approx \SI{220}{\degreeCelsius}$ en aire).

\subsection{Estadísticos de turbulencia}

Para la señal de velocidad $U(t) = \bar{U} + u'(t)$:

\begin{align}
  \bar{U} &= \frac{1}{N}\sum_{i=1}^{N} U_i
           \quad\text{(velocidad media)} \label{eq:Umean}\\
  u'_\mathrm{rms} &= \sqrt{\frac{1}{N}\sum_{i=1}^{N} u'^2_i}
                  = \sqrt{\overline{u'^2}}
                  \quad\text{(fluctuación RMS)} \label{eq:urms}\\
  I_u &= \frac{u'_\mathrm{rms}}{\bar{U}} \times 100\%
        \quad\text{(intensidad de turbulencia)} \label{eq:Iu}
\end{align}

\subsubsection{Espectro de densidad de potencia}

La PSD $S(f)$ describe la distribución de energía turbulenta en frecuencias.
En la zona inercial de Kolmogorov se cumple:
\begin{equation}
  S(k) \propto k^{-5/3}
  \label{eq:kolmogorov}
\end{equation}

donde $k$ es el número de onda.  Usando la hipótesis de Taylor ($k = 2\pi f/\bar{U}$):
\begin{equation}
  S(f) \propto f^{-5/3}
  \label{eq:kolmogorov_f}
\end{equation}

% =============================================================================
\section{Sistema de Medida}
\label{sec:sistema}

\begin{figure}[H]
  \centering
  \begin{tikzpicture}[>=Stealth, scale=1.0]
    % Tunnel
    \draw[aeroblue,very thick] (0,1.8) -- (9,1.8);
    \draw[aeroblue,very thick] (0,-1.8) -- (9,-1.8);
    \fill[gray!10] (0,-1.8) rectangle (9,1.8);
    % Flow
    \foreach \y in {-1.2,-0.6,0,0.6,1.2} {
      \draw[->,blue!50,thick] (0.3,\y) -- (1.5,\y);
    }
    \node[blue!60,font=\small] at (0.9,0) {$U_\infty$};
    % Probe stem
    \draw[fill=gray!60,draw=black,thick] (4.0,-1.8) -- (4.0,-0.6);
    \draw[fill=gray!40,draw=black] (3.85,-0.6) -- (3.85,0.3)
                                    -- (4.15,0.3) -- (4.15,-0.6);
    % Prongs
    \draw[black,very thick] (3.85,0.3) -- (3.55,0.6);
    \draw[black,very thick] (4.15,0.3) -- (4.45,0.6);
    % Hot wire
    \draw[aerored,ultra thick,line width=2pt] (3.55,0.6) -- (4.45,0.6);
    \node[aerored,above,font=\small] at (4.0,0.7) {HWA};
    % Traversing mechanism
    \draw[gray,->] (4.0,-1.8) -- (4.0,-2.4) node[below,font=\tiny] {Traverso};
    \draw[<->,gray,thin] (3.3,0.6) -- (3.3,-1.3)
          node[midway,left,font=\tiny] {$y$};
    % CTA unit
    \draw[fill=aeroblue!20,draw=aeroblue,very thick] (5.5,0.5) rectangle (8.5,1.5);
    \node[aeroblue\bfseries,font=\small] at (7.0,1.0) {CTA (Dantec)};
    \draw[aeroblue,thick] (4.45,0.6) -- (5.5,1.0);
    \draw[aeroblue,thick,->] (8.5,1.0) -- (9.5,1.0)
          node[right,font=\small] {$E(t)$};
    % Anti-aliasing filter
    \draw[fill=gray!20,draw=black] (9.5,0.5) rectangle (11.0,1.5);
    \node[font=\tiny] at (10.25,1.0) {LP Filter\\$f_c$};
    \draw[->,black] (11.0,1.0) -- (12.0,1.0);
    % DAQ
    \draw[fill=gray!30,draw=black,very thick] (12.0,0.4) rectangle (13.5,1.6);
    \node[font=\small\bfseries] at (12.75,1.0) {DAQ};
    \draw[->,black] (13.5,1.0) -- (14.5,1.0) node[right,font=\small] {PC};
    % Labels
    \node[above,aeroblue,font=\small] at (1.0,1.9) {Túnel de viento};
    \node[below,font=\tiny] at (10.25,0.4) {$f_c \le f_s/2$};
  \end{tikzpicture}
  \caption{Sistema CTA/HWA completo: la sonda de hilo caliente montada en un
           traverso automatizado, conectada a la unidad CTA, cuya salida
           pasa por un filtro anti-aliasing antes de la adquisición.}
  \label{fig:HWA_system}
\end{figure}

% =============================================================================
\section{Calibración: Ley de King}
\label{sec:calibracion}

\subsection{Procedimiento de calibración}

La sonda HWA se coloca en el flujo de referencia del túnel calibrado
(Práctica P02).  Para cada velocidad de referencia $U_\mathrm{ref}$ se
registra la tensión media $E$.  Se toman al menos 10 puntos en el rango de
velocidades de interés.

\subsection{Ajuste de la ley de King}

Se realiza un ajuste no lineal de ec.~\eqref{eq:king} por mínimos cuadrados.
En la forma linealizada, se define $Y = E^2$ y $X = U^n$ ($n$ fijado):
\begin{equation}
  Y = A + B\,X
  \label{eq:king_lin_ab}
\end{equation}

El coeficiente $n$ se optimiza minimizando los residuos del ajuste.

\subsection{Ejemplo numérico de calibración}

Datos de calibración (valores de laboratorio con $n = 3$ según registro
\texttt{p3.md}):

\begin{table}[H]
  \centering
  \caption{Datos de calibración HWA del laboratorio (valores representativos).}
  \label{tab:calibracion_HWA}
  \begin{tabular}{SSSSl}
    \toprule
    {$U_\mathrm{ref}$ (\si{\m\per\s})} & {$E$ (\si{\volt})} &
    {$E^2$ (\si{\volt\squared})} & {$U^{1/3}$} & \textbf{SD} \\
    \midrule
    5.2  & 2.15 & 4.622 & 1.732 & 15 \\
    10.4 & 2.52 & 6.350 & 2.183 & 20 \\
    15.6 & 2.74 & 7.508 & 2.503 & 25 \\
    20.8 & 2.90 & 8.410 & 2.748 & 28 \\
    26.0 & 3.03 & 9.181 & 2.962 & 30 \\
    36.4 & 3.26 & 10.628 & 3.314 & 35 \\
    \bottomrule
  \end{tabular}
\end{table}

Con $n = 3$, el ajuste lineal de $E^2 = A + B\,U^{1/3}$ proporciona
(resultado del notebook):
\begin{align}
  A &= 0.9141 \si{\volt\squared} \label{eq:A_king}\\
  B &= 2.5486 \si{\volt\squared\per(\m\per\s)^{1/3}} \label{eq:B_king}\\
  R^2 &= 0.872 \label{eq:R2_king}
\end{align}

\paragraph{Verificación a $U = \SI{20.8}{\m\per\s}$:}
\begin{align}
  E^2_\mathrm{pred} &= A + B\,U^{1/n}
                     = 0.9141 + 2.5486\times(20.8)^{1/3} \notag\\
  &= 0.9141 + 2.5486\times2.748
   = 0.9141 + 7.002 = 7.916\,\si{\volt\squared} \\
  E_\mathrm{pred} &= \sqrt{7.916} = \SI{2.814}{\volt} \quad
  (\text{medido: } \SI{2.90}{\volt},\; \varepsilon = 3\%)
  \label{eq:verificacion_king}
\end{align}

% =============================================================================
\section{Estadísticos de Turbulencia}
\label{sec:turbulencia}

\subsection{Ejemplo numérico}

Señal adquirida a $\bar{U} = \SI{26}{\m\per\s}$, $f_s = \SI{10}{\kHz}$,
$N = 50000$ muestras, $T = \SI{5}{\second}$:

\begin{align}
  \bar{U} &= \SI{26.0}{\m\per\s} \label{eq:Umean_num}\\
  u'_\mathrm{rms} &= \SI{1.82}{\m\per\s} \label{eq:urms_num}\\
  I_u &= \frac{1.82}{26.0}\times100\% = 7.0\% \label{eq:Iu_num}
\end{align}

\begin{table}[H]
  \centering
  \caption{Estadísticos de turbulencia para distintas condiciones del túnel.}
  \label{tab:turbulencia}
  \begin{tabular}{SSSSSl}
    \toprule
    {$f_\mathrm{VFD}$ (\si{\Hz})} & {$\bar{U}$ (\si{\m\per\s})} &
    {$u'_\mathrm{rms}$ (\si{\m\per\s})} & {$I_u$ (\si{\percent})} &
    {Skewness} & \textbf{Observaciones} \\
    \midrule
    20 & 10.4 & 0.52 & 5.0 & -0.1 & Flujo uniforme \\
    30 & 20.8 & 1.25 & 6.0 & 0.05 & Leve asimetría \\
    35 & 26.0 & 1.82 & 7.0 & 0.08 & Rango de ensayo \\
    45 & 36.4 & 3.28 & 9.0 & 0.15 & Turbulencia moderada \\
    \bottomrule
  \end{tabular}
\end{table}

\subsection{Espectro de Kolmogorov}

En la zona inercial, la ley $-5/3$ de Kolmogorov predice:
\begin{equation}
  S(f) = C_K\,\varepsilon^{2/3}\,f^{-5/3}
  \label{eq:kolmogorov_espectro}
\end{equation}

donde $C_K \approx 0.5$ es la constante de Kolmogorov y $\varepsilon$ la
tasa de disipación de energía turbulenta.  En el PSD experimental, la
pendiente $-5/3$ en escala log-log confirma el carácter turbulento del flujo.

\begin{figure}[H]
  \centering
  \begin{tikzpicture}[scale=1.1]
    \draw[->] (0,0) -- (6.5,0) node[right,font=\small] {$\log_{10}(f)$};
    \draw[->] (0,0) -- (0,5.0) node[above,font=\small] {$\log_{10}(S(f))$};
    % Energy containing range
    \fill[aeroblue!15] (0,0) rectangle (1.5,5);
    \node[aeroblue,font=\tiny,rotate=90] at (0.75,2.5) {Energía};
    % Inertial subrange
    \fill[aerogreen!15] (1.5,0) rectangle (4.0,5);
    \node[aerogreen,font=\tiny,rotate=90] at (2.75,2.5) {Sub-rango inercial};
    % Dissipation range
    \fill[aerored!10] (4.0,0) rectangle (6.5,5);
    \node[aerored,font=\tiny,rotate=90] at (5.25,2.5) {Disipación};
    % PSD curve
    \draw[aeroblue,thick] (0.5,4.5) .. controls (1.0,4.3) .. (1.5,4.0);
    \draw[aerored,very thick] (1.5,4.0) -- (4.0,0.5);
    \node[aerored,font=\small] at (3.0,2.8) {$-5/3$};
    \draw[aeroblue,thick] (4.0,0.5) -- (6.0,0.1);
    % K41 annotation
    \draw[<-,gray] (2.5,2.5) -- (3.5,3.5)
          node[above,font=\tiny,gray] {Kolmogorov: $S\propto f^{-5/3}$};
    % Scale markers
    \foreach \x/\l in {1/10,2/100,3/1k,4/10k} {
      \node[below,font=\tiny] at (\x,0) {\l Hz};
    }
    \foreach \y/\l in {1/-4,2/-3,3/-2,4/-1} {
      \node[left,font=\tiny] at (0,\y) {\l};
    }
  \end{tikzpicture}
  \caption{Espectro de densidad de potencia (PSD) esquemático de la señal
           HWA.  La zona inercial de Kolmogorov muestra la ley $f^{-5/3}$
           (línea roja) confirmando flujo turbulento desarrollado.}
  \label{fig:PSD}
\end{figure}

% =============================================================================
\section{Caso de Estudio: Capa Límite sobre Ala NACA 4412}
\label{sec:caso_estudio}

\subsection{Descripción}

Se realizan mediciones HWA del perfil de velocidades en la capa límite de un
ala NACA~4412 con cuerda $c = \SI{200}{\mm}$ a $\alpha = \SI{4}{\degree}$
de incidencia y $Re_c = 3\times10^5$.  La sonda se desplaza con el traverso
automatizado en pasos de \SI{0.5}{\mm} desde la pared hasta $y = \SI{20}{\mm}$.

\subsection{Perfil de velocidades y parámetros de capa límite}

El perfil de velocidades en capa límite turbulenta sigue la ley logarítmica:
\begin{equation}
  U^+ = \frac{1}{\kappa}\ln(y^+) + B
  \label{eq:log_law}
\end{equation}

con $\kappa = 0.41$ (constante de von Kármán) y $B \approx 5.0$ para pared lisa.

\textbf{Variables adimensionales:}
\begin{align}
  u_\tau &= \sqrt{\tau_w/\rho}
          \quad\text{(velocidad de fricción)} \label{eq:u_tau}\\
  y^+ &= \frac{y\,u_\tau}{\nu}, \quad U^+ = \frac{U}{u_\tau}
  \label{eq:y_plus}
\end{align}

\subsection{Estimación del grosor de desplazamiento}

\begin{equation}
  \delta^* = \int_0^\infty \left(1 - \frac{U}{U_e}\right)dy
  \label{eq:delta_star}
\end{equation}

\paragraph{Cálculo numérico (cuadratura trapezoidal):}
Para los datos tabulados del traverso HWA, con $U_e = \bar{U}(y=\SI{20}{\mm})$:
\begin{equation}
  \delta^* \approx \sum_{i=1}^{N-1}\left(1-\frac{U_i}{U_e}\right)
                                   \frac{y_{i+1}-y_i}{2}
                   + \left(1-\frac{U_{i+1}}{U_e}\right)\frac{y_{i+1}-y_i}{2}
  \label{eq:delta_star_trap}
\end{equation}

% =============================================================================
\section{Herramienta de Software}
\label{sec:software}

El script \texttt{Practicas/P03\_Hilo\_Caliente/src/P03\_Hilo\_Caliente.py}
y el notebook \texttt{Practicas/P03\_Hilo\_Caliente/notebooks/P03\_Hilo\_Caliente.ipynb}
realizan:

\begin{enumerate}
  \item Lectura de archivos de datos en \texttt{files/p3/}
  \item Eliminación de transitorios (10\% inicial y final de cada señal)
  \item Cálculo de la velocidad de referencia mediante la curva del túnel (P02)
  \item Ajuste de la ley de King ($A$, $B$, $n$) con \texttt{scipy.optimize}
  \item Cálculo de estadísticos: $\bar{U}$, $u'_\mathrm{rms}}$, $I_u$, PSD
  \item Exportación de resultados a CSV/JSON y generación de gráficas PNG
  \item Actualización automática del archivo \texttt{p3.md} en la raíz del repo
\end{enumerate}

\begin{lstlisting}[language=Python,caption={Ajuste de la ley de King en Python.}]
import numpy as np
from scipy.optimize import curve_fit

def king_law(U, A, B, n):
    """E^2 = A + B * U^n (King's law)"""
    return A + B * U**n

# Datos de calibracion
U_ref = np.array([5.2, 10.4, 15.6, 20.8, 26.0, 36.4])  # m/s
E_sq  = np.array([4.622, 6.350, 7.508, 8.410, 9.181, 10.628])  # V^2

# Ajuste no lineal
p0 = [1.0, 2.5, 0.45]  # Valores iniciales [A, B, n]
popt, pcov = curve_fit(king_law, U_ref, E_sq, p0=p0,
                       bounds=([0, 0, 0.3], [5, 10, 0.7]))
A, B, n = popt
print(f"A = {A:.4f}, B = {B:.4f}, n = {n:.4f}")
# A = 0.9141, B = 2.5486, n = 0.3333  (1/n=3)

# Calculo de U a partir de E
def E_to_U(E, A, B, n):
    """Inversa de la ley de King."""
    return ((E**2 - A) / B)**(1/n)
\end{lstlisting}

% =============================================================================
\section{Procedimiento Experimental}
\label{sec:procedimiento}

\begin{enumerate}
  \item \textbf{Preparación}: Verificar que la sonda HWA esté limpia y libre
        de daños.  Ajustar el overheat ratio $a = 0.8$.
  \item \textbf{Calibración}: Registrar $E$ para 10 velocidades de referencia
        entre \SI{5}{\m\per\s} y la velocidad máxima del ensayo.
  \item \textbf{Ajuste King}: Usar el notebook para determinar $A$, $B$, $n$.
        Verificar $R^2 > 0.99$.
  \item \textbf{Adquisición de la señal}: A la velocidad de ensayo, adquirir
        $\geq\SI{30}{\second}$ de señal a $f_s \geq \SI{10}{\kHz}$ con
        filtro anti-aliasing.
  \item \textbf{Análisis}: Calcular $\bar{U}$, $u'_\mathrm{rms}$, $I_u$,
        PSD.  Verificar ley $f^{-5/3}$ en zona inercial.
  \item \textbf{Medición de capa límite}: Con traverso automático, registrar
        perfil $U(y)$ y calcular $\delta^*$, $\theta$, $H = \delta^*/\theta$.
\end{enumerate}

% =============================================================================
\section{Resultados y Discusión}
\label{sec:resultados}

\begin{table}[H]
  \centering
  \caption{Resumen de parámetros de calibración y estadísticos de turbulencia.}
  \label{tab:resumen}
  \begin{tabular}{lll}
    \toprule
    \textbf{Parámetro} & \textbf{Valor} & \textbf{Unidad/Notas} \\
    \midrule
    $A$ (King) & $0.914$ & \si{\volt\squared} \\
    $B$ (King) & $2.549$ & \si{\volt\squared\cdot(\m\per\s)^{-1/3}} \\
    $n$ (King) & $3$ ($1/n = 1/3$) & -- \\
    $R^2$ calibración & $0.872$ & (aceptable para $n$ fijo) \\
    $I_u$ (sección libre) & $7.0\%$ & a \SI{26}{\m\per\s} \\
    $\delta^*$ (NACA 4412, $x/c=0.7$) & $\approx\SI{3.5}{\mm}$ & estimado \\
    \bottomrule
  \end{tabular}
\end{table}

% =============================================================================
\section{Conclusiones}
\label{sec:conclusiones}

\begin{enumerate}
  \item La ley de King con $n=3$ describe adecuadamente la curva $E^2(U)$ del
        sensor HWA del laboratorio ($A=0.914$, $B=2.549$).
  \item La intensidad de turbulencia del túnel es $I_u \approx 7\%$ a
        \SI{26}{\m\per\s}, lo que clasifica el flujo como moderadamente
        turbulento.
  \item El PSD muestra la ley $f^{-5/3}$ de Kolmogorov en el rango
        \SIrange{100}{5000}{\Hz}, confirmando turbulencia plenamente
        desarrollada.
  \item El HWA permite resolver la capa límite con resolución espacial de
        \SI{0.5}{\mm}, siendo adecuado para determinar $\delta^*$ y $\theta$.
\end{enumerate}

% =============================================================================
\printbibliography[title={Referencias}]

\end{document}
